% Latest publicly available official NeurIPS template as of 2026-03-23.
% A public NeurIPS 2026 style file was not yet available on neurips.cc,
% so this manuscript is aligned to the official NeurIPS 2025 template.
\documentclass{article}

\PassOptionsToPackage{numbers,compress}{natbib}
\usepackage[preprint]{neurips_2025}

\usepackage[utf8]{inputenc}
\usepackage[T1]{fontenc}
\usepackage{hyperref}
\usepackage{url}
\usepackage{booktabs}
\usepackage{amsfonts}
\usepackage{nicefrac}
\usepackage{microtype}
\usepackage{xcolor}
\usepackage{amsmath,amssymb}
\usepackage{graphicx}

\title{Activity Cliffs as Probes of Decision-Layer Collapse in Few-Shot Molecular Classification}
\author{Anonymous Authors}
\date{}

\begin{document}

\maketitle

\begin{abstract}
Average few-shot metrics can obscure failures in chemically critical local regions. In few-shot molecular
classification, a model may preserve relative ordering on hard pairs while still failing to form a usable decision
boundary in the presence of activity cliffs. We study this failure mode through \emph{CliffBench}, an activity-
cliff diagnostic benchmark and evaluation protocol instantiated on an assay-local few-shot molecular classification
substrate. CliffBench separates ranking-layer and decision-layer behavior through complementary metric families and
combines standard with adversarial episodic evaluation. On the rebuilt release, covering $6$ relaxed-profile tasks
and $2$ strict-profile tasks, we find that average-metric misalignment with cliff robustness is a supported trend,
while ranking--decision decoupling is strong enough to support a benchmark-level decision-layer collapse
interpretation. Among the full-strength baselines, ProtoNet is the strongest balanced model, whereas Random Forest
provides the clearest ranking-competent but decision-collapsed counterexample. A simple cliff-aware intervention
yields modest but consistent local gains over a vanilla $k$NN baseline, although not yet strong enough for a full
benchmark-level intervention claim. These results suggest that future progress in few-shot molecular learning will
require not only stronger representations, but also better local boundary formation and calibration in
cliff-sensitive regions.
\end{abstract}

\section{Introduction}

Machine learning is becoming an increasingly important component of molecular discovery, where predictive models
are used to prioritize compounds, guide hit expansion, and support lead optimization. In practice, however, many
molecular prediction tasks remain data-poor: each assay may provide only a small number of labeled examples, and
acquiring additional measurements is often expensive and slow. Few-shot learning (FSL) is therefore an appealing
paradigm for molecular property prediction, because it aims to transfer knowledge across tasks and adapt to a new
assay from only a small support set \cite{fsmol,protonet,maml}.

Strong average few-shot performance, however, does not by itself imply reliable decision making in the most
chemically important regions of structure--activity space. Medicinal chemistry is shaped by a tension between the
broad usefulness of the similarity principle and the sharp discontinuities that violate it. In many regions of a
structure--activity relationship (SAR) landscape, structurally similar molecules exhibit similar activity. In
\emph{activity cliffs}, by contrast, small structural changes can induce disproportionately large shifts in
potency or label \cite{activity_cliffs_review,moleculeace}. These cases are not peripheral curiosities. They are
precisely the regions in which a model must distinguish superficial similarity from decision-relevant local
structure.

This distinction becomes especially important in few-shot molecular classification. In a sparse episodic setting,
a model may learn enough relational structure to rank hard pairs correctly while still failing to place a usable
active/inactive boundary. We use the term \emph{decision-layer collapse} to describe this diagnostic pattern:
ranking behavior remains partly intact, but absolute decision behavior becomes fragile in cliff-heavy regions. In
other words, a model may still assign a higher score to the active molecule in a hard pair, yet fail to separate
the pair with a stable classification boundary. For downstream screening and prioritization, such behavior is
materially different from robust decision making.

Existing benchmarks do not make this failure mode easy to see. Standard few-shot molecular benchmarks emphasize
broad cross-task generalization and aggregate episode-level metrics \cite{fsmol}. Activity-cliff benchmarks, in
contrast, expose localized SAR discontinuities but are typically framed in conventional supervised settings rather
than assay-local episodic few-shot evaluation \cite{moleculeace}. As a result, the current evaluation landscape
leaves an important gap: we lack a benchmark substrate that directly tests whether few-shot molecular classifiers
can make stable decisions in the presence of high-similarity discordant pairs.

To address this gap, we introduce \emph{CliffBench}, an activity-cliff diagnostic benchmark and evaluation
protocol instantiated on an assay-local few-shot molecular classification substrate derived from FS-Mol. Rather
than treating cliffs as rare corner cases, CliffBench uses them as targeted probes of model reliability. The
protocol separates ranking-layer and decision-layer behavior through complementary metric families, combines
standard and adversarial episodic evaluation, and makes it possible to ask whether apparent few-shot competence
survives in the presence of chemically meaningful local discontinuities.

\subsection{Contributions}

This paper makes four contributions.

First, we motivate activity cliffs as high-value probes for evaluating few-shot molecular classifiers. Rather than
asking only whether a model performs well on average, we ask whether it remains reliable in local regions where
small structural changes have large decision consequences.

Second, we introduce \emph{CliffBench}, a benchmark and evaluation protocol that targets high-similarity
discordant and cliff-heavy regions in an assay-local few-shot setting. The framework combines standard and
adversarial episodes and explicitly separates ranking-layer from decision-layer evaluation.

Third, we show that average few-shot metrics and cliff-specific behavior do not move together cleanly. This
supports the view that global episode-level performance can mask localized failures, even when models look
competitive under standard aggregate summaries.

Fourth, we provide benchmark-level evidence consistent with a decision-layer collapse interpretation: some models
retain substantial hard-pair ranking ability while remaining weak on cliff decision and collapse-sensitive
metrics. We further show that a simple cliff-aware intervention yields modest but consistent local robustness
gains over a vanilla baseline, suggesting that localized supervision can improve cliff-side behavior even when the
overall effect remains limited.

More broadly, our goal is not to replace standard few-shot benchmarks, but to complement them with a diagnostic
view that is closer to the failure modes medicinal chemists care about. If few-shot molecular models are to be
useful in realistic discovery settings, they must do more than preserve coarse similarity-based ordering. They
must also remain reliable where small structural perturbations carry large decision consequences.

\section{Related Work}

\subsection{Few-Shot Molecular Property Prediction}

Few-shot learning has become a prominent strategy for molecular property prediction under assay-level data
scarcity. Existing work has explored metric-based, optimization-based, and transfer-based methods for adapting to
new molecular tasks from limited support examples \cite{fsmol,protonet,maml}. Benchmarks such
as FS-Mol were an important step forward because they standardized episodic evaluation across a large number of
assays and enabled systematic comparison among few-shot model families \cite{fsmol}. However, these benchmarks
were primarily designed to evaluate cross-task generalization at scale, rather than localized robustness in
chemically discontinuous regions.

\subsection{Activity Cliffs and SAR Discontinuities}

Activity cliffs are a longstanding topic in medicinal chemistry because they capture sharp violations of smooth
SAR assumptions: highly similar molecules can exhibit markedly different biological activity
\cite{activity_cliffs_review,moleculeace}. Recent benchmark efforts, most notably MoleculeACE, showed that overall
molecular prediction quality and cliff-specific performance can diverge substantially, and that models with strong
global metrics may still perform poorly on cliff subsets \cite{moleculeace}. This line of work established the
importance of cliff-aware evaluation, but it has largely focused on standard supervised learning rather than
assay-local few-shot episodes. Our work builds on this insight by bringing cliff-centered evaluation into a few-
shot setting.

\subsection{Shortcut Learning, Calibration, and Boundary Fragility}

A growing literature argues that molecular learning systems can rely on shortcut-like signals, including global
similarity or scaffold-level regularities, instead of learning the localized determinants that drive true
functional differences \cite{shortcut_learning,molecular_shortcuts}. This concern is particularly relevant for
activity cliffs, where scaffold-level similarity can be high while decision-relevant local changes remain subtle.
More broadly, recent work in machine learning and molecular foundation models has highlighted objective--metric
misalignment and reliability limits: models can optimize ranking-oriented or contrastive objectives while remaining
poorly calibrated for absolute decisions \cite{metric_misalignment,uncertainty_mol}. Our work is aligned with this
perspective, but focuses on an assay-local few-shot regime in which sparse support sets can make decision
boundaries especially fragile.

\subsection{Positioning of This Work}

Our setting sits at the intersection of these three lines of research. Relative to standard few-shot molecular
benchmarks, we focus on localized high-similarity discordant regions rather than only global episodic averages.
Relative to activity-cliff benchmarks, we operate in a true few-shot episodic regime. Relative to shortcut-
learning and calibration analyses, we contribute a benchmark-level diagnostic protocol rather than a purely
mechanistic claim. The result is a framework for evaluating whether a few-shot molecular classifier can do more
than rank broadly plausible candidates: it can be tested on whether it maintains a usable decision boundary where
SAR is most discontinuous.

\section{Problem Setting}

\subsection{Assay-Local Few-Shot Molecular Classification}

We consider assay-level binary few-shot molecular classification. For an assay $t$, let $V_t$ denote the set of
valid molecules, with binary labels $y_i \in \{0,1\}$ and assay-local continuous activity values $r_i \in
\mathbb{R}$ for each molecule $m_i \in V_t$. We write $V_t^+$ and $V_t^-$ for the active and inactive subsets,
respectively.

Evaluation is episodic. In each episode, a model receives a small labeled support set and must predict labels or
scores for a query set from the same assay. The few-shot goal is not merely to rank molecules coarsely by
plausibility, but to transfer across assays and form a useful decision rule from limited support examples.

This setting is particularly challenging in regions where structural similarity is high but label discordance
persists. Such cases are common in medicinal chemistry and are especially informative for evaluating whether a
model has learned a decision-relevant representation rather than a broad similarity heuristic.

\subsection{High-Similarity Discordant Pairs, Activity Cliffs, and Controls}

For each assay, we compute molecular similarity using Tanimoto similarity over Morgan fingerprints. Let
$\mathrm{Sim}(m_i,m_j)$ denote the similarity between molecules $m_i$ and $m_j$. Following an assay-local
definition, we first identify the set of high-similarity discordant pairs
\[
H_t^{\mathrm{disc}} = \{(i,j)\mid i \in V_t^+,\, j \in V_t^-,\, \mathrm{Sim}(m_i,m_j) \ge \tau \}.
\]

We then partition this set into cliff and non-cliff controls using the assay-local activity gap:
\[
C_t = \{(i,j) \in H_t^{\mathrm{disc}} \mid |r_i-r_j| \ge \delta \},
\]
\[
D_t = \{(i,j) \in H_t^{\mathrm{disc}} \mid |r_i-r_j| < \delta \}.
\]

Here, $C_t$ contains activity-cliff pairs, while $D_t$ serves as a high-similarity non-cliff control set.
Throughout the benchmark, pair direction is fixed as active $\rightarrow$ inactive. We also define same-scaffold
subsets by restricting to pairs that share the same non-empty Bemis--Murcko scaffold:
\[
C_t^{\mathrm{scaf}} \subseteq C_t, \qquad D_t^{\mathrm{scaf}} \subseteq D_t.
\]

This construction separates three related but distinct notions: broad assay-level classification, hard high-
similarity discordant ranking, and true cliff-side decision making.

\subsection{Diagnostic View: Ranking Layer vs. Decision Layer}

Let a model output an active score $s(x) \in \mathbb{R}$ for molecule $x$. For pairwise ranking, we define the
pairwise success function
\[
\phi(i,j)=
\begin{cases}
1, & s(m_i) > s(m_j),\\
0.5, & s(m_i) = s(m_j),\\
0, & s(m_i) < s(m_j).
\end{cases}
\]
Because pair direction is fixed as active $\rightarrow$ inactive, higher values indicate better ranking of the
active molecule above the inactive one.

We use this to define ranking-layer metrics. For query-side cliff pairs,
\[
\mathrm{Q\mbox{-}PSR} = \frac{1}{|Q_t^{\mathrm{cliff}}|} \sum_{(i,j)\in Q_t^{\mathrm{cliff}}} \phi(i,j),
\]
where $Q_t^{\mathrm{cliff}}$ denotes the set of cliff pairs induced by the query set in an episode. For
adversarially injected support-query cliff pairs, we analogously compute $\mathrm{SQ\mbox{-}PSR}$. Same-scaffold
variants restrict the averaging set to pairs in $C_t^{\mathrm{scaf}}$.

To evaluate decision behavior, we convert scores to discrete predictions $\hat{y}_i \in \{0,1\}$ using the model's
fixed reporting rule. We then compute cliff-side balanced accuracy on cliff-associated query molecules:
\[
\mathrm{C\mbox{-}BAcc} = \frac{1}{2}(\mathrm{TPR}_{\mathrm{cliff}} + \mathrm{TNR}_{\mathrm{cliff}}).
\]

We further define the success-but-collapse rate through the collapse statistic
\[
\mathrm{SCR} = \frac{1}{|P_t|} \sum_{(i,j)\in P_t} \mathbb{I}\!\left[\hat{y}_i=\hat{y}_j\right],
\]
where $P_t$ denotes the relevant high-similarity discordant pair set for the evaluation slice. Intuitively,
$\mathrm{SCR}$ measures how often a model assigns the same discrete class to both molecules in a discordant pair;
lower values are better. Same-scaffold collapse is measured by $\mathrm{SS\mbox{-}SCR}$.

These metrics separate two behaviors that are often conflated in standard evaluation. A model may retain
substantial ranking competence on hard pairs while still failing to establish a usable classification boundary. We
refer to this ranking--decision decoupling as \emph{decision-layer collapse}.

\subsection{Research Questions}

Our evaluation is organized around three questions.

\paragraph{H1: Metric Misalignment.}
Do average few-shot metrics and cliff-specific metrics move together cleanly across models, or can strong average
performance mask weak cliff robustness?

\paragraph{H2: Decision-Layer Collapse.}
Can a model preserve ranking behavior on hard high-similarity pairs while remaining weak on cliff decision and
collapse-sensitive metrics?

\paragraph{H3: Cliff-Aware Mitigation.}
Does introducing cliff-aware supervision or episode construction improve cliff-side behavior and reduce collapse
without substantially harming non-cliff controls?

\section{CliffBench: Benchmark Construction and Evaluation Protocol}

\subsection{Assay-Local Benchmark Substrate}

We instantiate \emph{CliffBench} on an assay-local few-shot molecular classification substrate derived from FS-Mol
\cite{fsmol}. All preprocessing, pair mining, episode construction, and metric computation are
performed within an assay. We do not mix continuous activities, pair definitions, or scaffold logic across assays.

A molecule is retained only if it passes a fixed legality filter, including finite assay-local activity, precise
relation type, valid RDKit parsing and sanitization, and canonical isomeric identity handling. Duplicate records
are resolved assay-locally using deterministic rules. This design ensures that cliff and non-cliff events remain
chemically meaningful within each task rather than becoming artifacts of cross-assay aggregation.

\subsection{Profiles and Task Eligibility}

CliffBench uses two benchmark profiles. The \emph{relaxed} profile is the primary benchmark, while the
\emph{strict} profile is retained as a sharper mechanism stress test. In the rebuilt release, the profiles use
\[
(\tau,\delta) = (0.80,1.0) \quad \text{for relaxed},
\]
\[
(\tau,\delta) = (0.85,1.0) \quad \text{for strict},
\]
with common minimum requirements on cliff and non-cliff support.

An assay is benchmark-eligible only if it satisfies minimum constraints on valid molecule count, class balance,
high-similarity discordant support, cliff support, non-cliff control support, and anchor-side coverage.
Adversarial eligibility further requires sufficient support-query matching capacity, quantified through the
maximum matching size $M_{\mathrm{avail}}(t)$ between active anchors and cliff negatives.

This dual-profile design serves two complementary purposes. The relaxed profile increases assay coverage and
functions as the main comparison substrate. The strict profile concentrates the cliff-heavy regime and is better
suited for mechanism inspection and stress testing.

\subsection{Pair Mining and Hard Negative Construction}

For each assay, we mine high-similarity discordant pairs from active--inactive molecule pairs satisfying the
profile-specific similarity threshold. Activity cliffs and non-cliff controls are then separated by the assay-
local activity gap threshold. We additionally identify same-scaffold subsets using Bemis--Murcko scaffolds.

To support adversarial episode construction, we build an anchor-specific hard negative pool. For each active
anchor, candidate inactive partners are ranked by decreasing similarity, then by decreasing activity gap, with
deterministic tie-breaking. This provides a consistent source of chemically close but label-opposed negatives for
adversarial injection.

\subsection{Standard and Adversarial Episodes}

CliffBench uses two episode types.

\paragraph{Standard episodes.}
Standard episodes follow the frozen few-shot configuration of the underlying benchmark substrate. They provide the
baseline estimate of assay-level few-shot behavior under a conventional episodic split.

\paragraph{Adversarial episodes.}
Adversarial episodes explicitly inject support-query cliff structure. For each assay, we construct a bipartite
graph between active anchors and cliff negatives, and compute the maximum matching size $M_{\mathrm{avail}}(t)$.
The injection count is then
\[
m = \min\!\left(\left\lfloor 0.5|Q^-| \right\rfloor,\ |S^+|,\ M_{\mathrm{avail}}(t)\right),
\]
where $Q^-$ denotes the negative query set and $S^+$ the positive support set. If $m<2$, the assay does not
contribute to adversarial metrics.

This design makes adversarial evaluation targeted rather than random. Instead of hoping that cliff events appear
naturally in the query set, CliffBench guarantees that the model is tested on structurally similar but decision-
critical positive--negative relationships.

\subsection{Metric Families and Reporting}

CliffBench reports three complementary metric families.

\paragraph{Ranking-layer metrics.}
These include $\mathrm{Q\mbox{-}PSR}$, $\mathrm{SQ\mbox{-}PSR}$, and their same-scaffold counterparts. They
measure whether the model places the active molecule above the inactive molecule in hard pair comparisons.

\paragraph{Decision-layer metrics.}
These include $\mathrm{C\mbox{-}BAcc}$, $\mathrm{NC\mbox{-}BAcc}$, $\mathrm{SCR}$, and $\mathrm{SS\mbox{-}SCR}$.
They measure whether the model forms a usable boundary on cliff-side molecules and whether discordant pairs
collapse to identical discrete predictions.

\paragraph{Average metrics.}
To preserve comparability with standard few-shot reporting, we also report average precision and $\Delta$AUPRC,
where $\Delta$AUPRC subtracts the positive-class base rate in the query set.

Together, these metrics expose a failure mode that standard aggregate reporting can hide. A model can look
competitive under average metrics or even under pairwise ranking metrics, while remaining weak on cliff-specific
decision metrics. This separation is central to our diagnostic view.

\section{Experimental Setup}

\subsection{Benchmark Coverage}

We evaluate on the rebuilt CliffBench release. Across $157$ raw test assays considered during rebuilding, the
current release contains $6$ relaxed-eligible assays and $2$ strict-eligible assays.
Following the benchmark positioning, we use \emph{relaxed} as the primary comparison substrate and \emph{strict}
as supporting mechanism context.

This distinction matters for claim strength. The relaxed profile provides the most defensible basis for benchmark-
level comparisons, while the strict profile is best interpreted as a stress-test slice with lower coverage.

\subsection{Model Suite}

Our primary full-strength comparison set consists of four model families:
\begin{itemize}
    \item \textbf{kNN}: a local metric baseline;
    \item \textbf{RF}: a strong classical supervised baseline;
    \item \textbf{ProtoNet}: a metric-based few-shot baseline with support-side scoring;
    \item \textbf{kNN-cliff-aware}: a localized cliff-aware intervention baseline.
\end{itemize}

We also retain \textbf{MAML} as exploratory compatibility context. However, in the current rebuilt release its
rows are generated through a legacy compatibility path with only $3$ episodes per task/seed/split, rather than the
full evaluation budget. We therefore do not treat MAML as part of the strongest final claim substrate.

\subsection{Episode Configuration}

Unless otherwise noted, episodes use a $2$-way few-shot setting with $16$ support molecules per class and $16$
query molecules per class. For each task and seed, we evaluate
\[
400 \text{ standard episodes} + 400 \text{ adversarial episodes}.
\]
We use seeds $0,\dots,4$ throughout the main evaluation.

This configuration is intentionally large enough to stabilize episode-level summaries while remaining compatible
with assay-local coverage constraints in the rebuilt release.

\subsection{Prediction Rules and Support-Side Scoring}

Each model outputs an active score $s(x)$. Ranking-layer metrics operate directly on these continuous scores.
Decision-layer metrics require discrete predictions derived from the model's fixed reporting rule. For methods
that report support-query ranking metrics, we require an explicit support-side scoring mechanism, such as
prototype-based support scoring or post-adaptation classifier outputs, so that support-query comparisons are
operationally well-defined.

This is important because a model should not be credited with support-query ranking performance if it only
produces query-side scores.

\subsection{Aggregation and Statistical Comparison}

Evaluation proceeds in two stages. First, within each task, seed, split, and metric, we average over valid
episodes only. For each metric we also report coverage, the number of valid episodes, and the mean number of valid
pairs per valid episode. Second, we macro-average task-level summaries across tasks and report $95\%$ confidence
intervals using task-level bootstrap resampling with $10{,}000$ iterations.

For model comparisons, we use task-level paired bootstrap deltas. This design keeps the comparison aligned to
assay-level variation rather than inflating confidence through raw episode counts. Because assay coverage is still
limited in the rebuilt release, we interpret confidence intervals conservatively and distinguish between
\emph{formal claim}, \emph{supported trend}, and \emph{exploratory} evidence tiers.

\subsection{Scope of Main Claims}

Our strongest claim wording is based on the relaxed-profile full-strength model set: \texttt{kNN}, \texttt{RF},
\texttt{ProtoNet}, and \texttt{kNN-cliff-aware}. Under this scope, the current release supports a formal
benchmark-level interpretation of decision-layer collapse, while the metric-misalignment and intervention stories
are better stated as supported trends rather than final claims.

We adopt this claim discipline throughout the paper. The goal is to characterize the current evidence boundary
clearly, rather than to overstate what can be concluded from a rebuilt benchmark with limited assay coverage.

\section{Results}

Unless otherwise noted, all quantitative results in this section refer to the relaxed-profile full-strength model
set: \texttt{kNN}, \texttt{RF}, \texttt{ProtoNet}, and \texttt{kNN-cliff-aware}. We use this substrate for main-
text interpretation because it is the current strongest basis for benchmark-level comparison.

\subsection{Aggregate Snapshot}

Table~\ref{tab:relaxed-main} summarizes the main relaxed-profile aggregate metrics. Two patterns are immediately
visible. First, ProtoNet is the strongest balanced full-episode baseline across both the standard and adversarial
slices. Second, RF is the clearest example of a model that remains highly competitive at hard-pair ranking while
still looking fragile on cliff-side decision behavior.

\begin{table}[t]
\centering
\caption{Aggregate relaxed-profile results for the full-strength model set. Higher is better for
$\Delta$AUPRC, C-BAcc, Q-PSR, and SQ-PSR; lower is better for SCR.}
\label{tab:relaxed-main}
\resizebox{\textwidth}{!}{%
\begin{tabular}{lcccccccc}
\toprule
Model & std $\Delta$AUPRC & std C-BAcc & std Q-PSR & std SCR & adv $\Delta$AUPRC & adv C-BAcc & adv SQ-PSR & adv SCR \\
\midrule
\texttt{kNN} & 0.105021 & 0.508397 & 0.513157 & 0.925098 & 0.015338 & 0.474543 & 0.528708 & 0.939580 \\
\texttt{RF} & 0.187063 & 0.514127 & 0.591827 & 0.899767 & -0.009912 & 0.510866 & 0.910745 & 0.933810 \\
\texttt{ProtoNet} & 0.203395 & 0.544073 & 0.643762 & 0.804123 & 0.057186 & 0.547411 & 0.778281 & 0.840986 \\
\texttt{kNN-cliff-aware} & 0.109866 & 0.514785 & 0.518979 & 0.897206 & 0.030065 & 0.489892 & 0.539401 & 0.918337 \\
\bottomrule
\end{tabular}%
}
\end{table}

\subsection{Average Metrics Do Not Cleanly Track Cliff Robustness}

Our first question is whether standard few-shot averages provide a faithful ordering of cliff robustness. The
answer is only partially yes. ProtoNet is the strongest model both on average-oriented metrics and on several
cliff-sensitive metrics, but the broader model ordering does not collapse into a single one-dimensional ranking. In
particular, the transition from \texttt{kNN} to \texttt{RF} shows the clearest masking pattern: on the standard
split, paired comparisons show gains in $\Delta$AUPRC and Q-PSR, but the direct gain in cliff-side balanced
accuracy is small and its interval does not cleanly separate from zero.

This is why we treat the metric-misalignment story as a supported trend rather than a final claim. The signal is
real, but the current rebuilt release still has only six relaxed-profile tasks, and the cliff-versus-control gap
is not directionally stable enough across the full model set to justify stronger wording. In other words, average
few-shot metrics can mask localized cliff failures, but the present artifact supports this as a credible trend
rather than a fully settled benchmark-level conclusion.

\subsection{Activity Cliffs Expose Ranking--Decision Decoupling}

The strongest result in the rebuilt release is the ranking--decision split itself. RF provides the sharpest
counterexample. On the adversarial split, it reaches $\mathrm{SQ\mbox{-}PSR}=0.910745$, indicating extremely
strong support-query hard-pair ranking, yet its cliff-side decision behavior remains weak, with
$\mathrm{C\mbox{-}BAcc}=0.510866$ and $\mathrm{SCR}=0.933810$. This is precisely the pattern we mean by
decision-layer collapse: the model can often order hard pairs correctly while still failing to establish a usable
decision boundary.

This is not an outlier-driven story. At the task level, RF beats \texttt{kNN} on adversarial SQ-PSR in all
$6/6$ relaxed tasks, but beats it on adversarial C-BAcc in only $3/6$ tasks. The ranking signal is broad and
stable; the decision signal is not. ProtoNet shows the same split in a milder form, but remains substantially more
balanced than RF. It is stronger than \texttt{kNN} on adversarial C-BAcc in $6/6$ tasks and on adversarial SCR in
$5/6$ tasks, which is why it is best interpreted as the strongest balanced baseline rather than as a model that
eliminates collapse entirely.

Same-scaffold results reinforce this picture. Ranking remains harder when evaluation is restricted to the same-
scaffold slice, with RF dropping from adversarial SQ-PSR $0.910745$ to adversarial SS-SQ-PSR $0.897061$, and
ProtoNet dropping from $0.778281$ to $0.763198$. This same-scaffold brittleness is consistent with the broader
interpretation that localized cliff decisions remain difficult even when global similarity signals are strong.

\subsection{Cliff-Aware Intervention Yields Modest but Consistent Local Gains}

Our third question asks whether a localized cliff-aware intervention can improve cliff-side behavior without
substantially harming non-cliff controls. The answer is encouraging but modest. Relative to vanilla \texttt{kNN},
\texttt{kNN-cliff-aware} improves standard C-BAcc from $0.508397$ to $0.514785$ and reduces standard SCR from
$0.925098$ to $0.897206$. In paired comparisons, the standard C-BAcc gain is $+0.006388$ with a $95\%$ interval
$[0.002617, 0.010193]$, while standard SCR improves by $-0.027892$ with interval
$[-0.056174, -0.002390]$.

The intervention also helps under adversarial evaluation, although again with limited magnitude. Raw adversarial
C-BAcc rises from $0.474543$ to $0.489892$, and adversarial SCR drops from $0.939580$ to $0.918337$. The paired
adversarial SCR improvement is $-0.021242$ with interval $[-0.045223, -0.004900]$. At the task level, the
intervention improves standard C-BAcc in $5/6$ tasks and adversarial SCR in $5/6$ tasks, indicating that the
effect is not driven by only one or two assays.

At the same time, the intervention does not yet support a strong benchmark-level mitigation claim. Its adversarial
SQ-PSR direction splits $3/3$ across tasks, and the paired adversarial C-BAcc interval does not cleanly separate
from zero. For this reason, we interpret the intervention result as a supported trend: localized supervision can
improve cliff-side decision quality and reduce collapse, but the current rebuilt release does not yet justify a
stronger claim.

\subsection{Overall Model Profiles}

Taken together, the rebuilt relaxed benchmark supports three distinct model profiles. ProtoNet is the strongest
balanced baseline, combining the best average metrics with the most convincing cliff-side decision behavior in the
full-strength model set. RF is the clearest ranking-competent but decision-collapsed baseline, making it the most
useful counterexample for the decision-layer collapse interpretation. The cliff-aware \texttt{kNN} variant acts as
a local rectifier: it does not dominate the benchmark, but it consistently improves several cliff and collapse
metrics over the vanilla \texttt{kNN} baseline. These profiles are more informative than a single global ranking of
models because they expose which aspect of few-shot molecular competence each method does, or does not, provide.

\section{Discussion and Limitations}

Our results support a simple but practically important view of few-shot molecular reliability. In cliff-heavy
regions, it is not enough for a model to preserve coarse similarity-aware ordering. What matters is whether that
ordering can be converted into a usable local decision rule. CliffBench makes this distinction visible by
separating ranking-layer and decision-layer behavior. Under this view, the strongest message of the rebuilt release
is not merely that some models are better than others, but that different models fail in different ways. RF is
useful precisely because it reveals how strong hard-pair ranking can coexist with weak cliff-side decisions,
whereas ProtoNet is useful because it shows that a more balanced few-shot baseline is possible without eliminating
the underlying difficulty.

These findings also suggest a methodological shift for future molecular few-shot work. Progress should not be
measured only by better aggregate averages or stronger pairwise ranking. If activity cliffs are treated as
decision-critical probes, then local boundary formation, calibration, and robustness under support-query
perturbation become central design targets. The modest but consistent gains of the cliff-aware \texttt{kNN}
intervention point in this direction. They do not solve the problem, but they indicate that explicit attention to
cliff structure can improve local decision quality without obviously destroying non-cliff behavior.

At the same time, the current study has several limitations. First, the rebuilt benchmark remains narrow. The main
comparison substrate contains only $6$ relaxed-profile tasks and $2$ strict-profile tasks, which is sufficient for
a meaningful diagnostic study but still limits the stability of model-ordering and intervention claims. This is
the main reason why we treat H1 and H3 as supported trends rather than final benchmark-level claims. Second, our
evidence is benchmark-level and behavioral. While the observed patterns are consistent with shortcut-like reliance
on global similarity and with local boundary fragility, they do not by themselves constitute causal proof of a
particular representation or geometric mechanism. Third, the current full-strength comparison set is intentionally
restricted. MAML is retained only as exploratory compatibility context because its rebuilt release path does not
yet match the full evaluation budget of the other models. Finally, although the release already exposes both
standard and adversarial cliff-sensitive metrics, the present paper does not yet include direct threshold-fragility
analysis or richer control-only artifact reporting; both would make the decision-layer interpretation more explicit.

These limitations define a clear path forward. The next stage is not to make the narrative stronger by language,
but to make it stronger by evidence: larger task coverage, cleaner control-gap reporting, explicit threshold
analysis, and broader baseline coverage would all sharpen the present conclusions. Even under the current release
boundary, however, the central diagnostic point is already clear: few-shot molecular models should be judged not
only by whether they can rank plausible candidates, but also by whether they can sustain reliable decisions in the
regions where medicinal chemistry is most discontinuous.

\section{Conclusion}

This paper introduced CliffBench, an activity-cliff diagnostic benchmark and evaluation protocol for assay-local
few-shot molecular classification. By separating ranking-layer from decision-layer evaluation, CliffBench exposes a
failure mode that standard aggregate reporting can hide: models may preserve hard-pair ordering while remaining
fragile at the level of cliff-side decisions. On the rebuilt release, this ranking--decision decoupling is strong
enough to support a benchmark-level decision-layer collapse interpretation.

Our empirical picture is therefore deliberately asymmetric. Average-metric misalignment with cliff robustness is
credible but not yet stable enough for a strongest-form claim. The cliff-aware intervention signal is real but
still modest. The clearest current result is instead that activity cliffs function as high-value probes of whether
a few-shot molecular model has learned a usable local decision boundary. In this sense, CliffBench is not meant to
replace standard few-shot benchmarks, but to complement them with a reliability-oriented diagnostic layer.

If few-shot molecular learning is to become useful in realistic discovery settings, it must do more than rank
similar molecules in roughly the right order. It must also remain trustworthy where small structural changes carry
large decision consequences. Activity cliffs provide a principled way to test that requirement, and the present
results show that this requirement is not yet routinely met.

\bibliographystyle{unsrtnat}
\bibliography{references}

\end{document}
